\documentclass[9pt]{article}
\usepackage[a4paper,top=8mm,bottom=8mm,left=12mm,right=12mm]{geometry}
\usepackage[T1]{fontenc}
\usepackage[utf8]{inputenc}
\usepackage[spanish,es-nodecimaldot]{babel}
\usepackage[scaled=0.95]{helvet}
\renewcommand{\familydefault}{\sfdefault}
\usepackage{sansmath}
\sansmath
\usepackage{amsmath,amssymb}
\usepackage{tikz}
\usetikzlibrary{arrows.meta,calc,positioning}
\usepackage{tabularx,array}
\usepackage{enumitem}
\usepackage{ragged2e}
\usepackage{microtype}
\usepackage{xcolor}
\pagestyle{empty}
\setlength{\parindent}{0pt}
\setlength{\parskip}{0pt}
\setlength{\tabcolsep}{2pt}
\setlist[itemize]{leftmargin=3.4mm,itemsep=0.3mm,topsep=0mm,parsep=0mm,partopsep=0mm,label=\raisebox{0.2ex}{\scriptsize$\blacksquare$}}
\newcommand{\problem}[1]{\textbf{Problema #1:}}
\newcommand{\vect}[1]{\vec{#1}}
\newcommand{\tightline}{\vspace{0.55mm}\hrule height 0.45pt\vspace{0.75mm}}
\begin{document}
\fontsize{9.65}{12.45}\selectfont

% Encabezado
\begin{minipage}[c]{0.18\textwidth}
\centering
\begin{tikzpicture}[x=1cm,y=1cm]
  \draw[line width=0.45pt] (0,0) circle (0.66);
  \node[font=\bfseries\fontsize{8.2}{8.3}\selectfont] at (0,0.16) {USACH};
  \node[font=\bfseries\fontsize{3.25}{3.45}\selectfont,align=center] at (0,-0.12) {DEPARTAMENTO\\DE FÍSICA};
\end{tikzpicture}
\end{minipage}%
\begin{minipage}[c]{0.48\textwidth}
\centering
{\bfseries\fontsize{13.2}{13.5}\selectfont PEP 1}\par
{\fontsize{10.0}{10.4}\selectfont Forma B}\par
{\fontsize{10.0}{10.4}\selectfont Tiempo: 80 min.}
\end{minipage}%
\begin{minipage}[c]{0.34\textwidth}
\raggedleft
{\bfseries\fontsize{9.4}{9.8}\selectfont Física I para Ingeniería}\par
{\bfseries\fontsize{9.2}{9.6}\selectfont Cód: 10103-10140}\par
{\bfseries\fontsize{9.2}{9.6}\selectfont 1er Semestre, 2026}
\end{minipage}

\vspace{1.2mm}

% Caja de indicaciones
\setlength{\fboxsep}{1.4mm}
\noindent\fbox{%
\begin{minipage}{0.976\textwidth}
\begin{tabularx}{\textwidth}{@{}X@{\hspace{5mm}}X@{}}
\begin{minipage}[t]{\linewidth}
\textbf{Indicaciones:}
\begin{itemize}
  \item Tu procedimiento debe ser \textbf{CLARO Y ORDENADO}. Identifica claramente a qué inciso y problema corresponde cada desarrollo.
  \item Desarrollos con lápiz grafito no serán re-corregidos.
  \item Respuestas numéricas deben ser encerradas en un recuadro y presentar unidades de medida en S.I. (de no estar presente no podrá obtener el puntaje completo del problema).
\end{itemize}
\end{minipage}
&
\begin{minipage}[t]{\linewidth}
\vspace{2.2mm}
\begin{itemize}
  \item Considera $g=9{,}8\,\mathrm{m/s^2}$, en el caso que sea necesario utilizarlo.
  \item Sólo se permite uso de calculadora personal no programable.
  \item Solamente aproxima tu \textbf{RESULTADO FINAL} y utiliza 2 decimales para ello.
\end{itemize}
\end{minipage}
\end{tabularx}
\vspace{0.55mm}\hrule height 0.35pt\vspace{0.7mm}
\centering\bfseries\fontsize{8.4}{9.3}\selectfont
Sólo se permiten preguntas relacionadas exclusivamente con la redacción del texto. Las personas a cargo de la sala durante la evaluación NO tienen permitido responder consultas acerca del contenido que se evalúa y la Coordinación no asumirá responsabilidad por cualquier orientación o comentario emitido.
\end{minipage}}

\vspace{1.15mm}

% Problema 1 y gráfico
\begin{minipage}[t]{0.59\textwidth}
\justifying
\problem{1} Se hacen pruebas para la implementación de un camión de reparto autónomo de masa $M=1052\,\mathrm{kg}$ y su batería tiene una autonomía de 6 horas. Este utiliza cámaras con inteligencia artificial que le indican cuando acelerar o frenar al acercarse a estaciones de entrega o retiro de paquetes. Como primer ensayo, el camión se desplaza en línea recta entre dos estaciones. El GPS del camión presenta siguiente gráfico $v$ vs. $t$, que describe su movimiento completo en tres etapas. Parte del ensayo es simular una situación alta demanda y se manda un segundo vehículo autónomo de apoyo de masa $m=2510\,\mathrm{kg}$ que parte desde el mismo origen con un retraso de 2 segundos, a velocidad inicial $v_0$ desconocida y aceleración constante $\vect{a}=(3{,}0\,\mathrm{m/s^2})\,\hat{\imath}$. Ambos se encuentran una vez en la segunda etapa y el vehículo autónomo de apoyo alcanza una rapidez de $v_f=26\,\mathrm{m/s}$.
\end{minipage}\hfill
\begin{minipage}[t]{0.385\textwidth}
\centering
\vspace{-0.5mm}
\begin{tikzpicture}[x=0.235cm,y=0.105cm]
  \draw[-{Latex[length=2mm]},line width=0.58pt] (0,0) -- (28.4,0);
  \node[above left=1pt,font=\fontsize{8.2}{8.2}\selectfont] at (28.0,0) {$t\,(\mathrm{s})$};
  \draw[-{Latex[length=2mm]},line width=0.58pt] (0,0) -- (0,31) node[right=2pt,font=\fontsize{8.2}{8.2}\selectfont] {$v\,(\mathrm{m/s})$};
  \draw[dashed,line width=0.35pt] (0,24) -- (28,24);
  \draw[dashed,line width=0.35pt] (8,0) -- (8,30);
  \draw[dashed,line width=0.35pt] (20,0) -- (20,30);
  \draw[line width=0.78pt] (0,0) -- (8,24) -- (20,24) -- (26,0);
  \node[below left=1pt,font=\fontsize{7.5}{7.5}\selectfont] at (0,0) {$0$};
  \node[below=2pt,font=\fontsize{7.5}{7.5}\selectfont] at (8,0) {$8$};
  \node[below=2pt,font=\fontsize{7.5}{7.5}\selectfont] at (20,0) {$20$};
  \node[below=2pt,font=\fontsize{7.5}{7.5}\selectfont] at (26,0) {$26$};
  \node[left=2pt,font=\fontsize{7.5}{7.5}\selectfont] at (0,24) {$24$};
\end{tikzpicture}
\vspace{-0.7mm}

{\fontsize{8.2}{8.7}\selectfont Figura 1: Gráfica velocidad-tiempo.}
\end{minipage}

\vspace{0.55mm}
\begin{tabularx}{\textwidth}{@{}>{\raggedright\arraybackslash}X>{\raggedleft\arraybackslash}p{25mm}@{}}
(a) Determina la aceleración del camión en la primera y en la tercera etapa, $(a_1)$ y $(a_3)$, respectivamente. & (5 puntos)\\[-0.15mm]
(b) Utilizando los resultados anteriores, expresa las ecuaciones de posición $x(t)$ para cada una de esas etapas, tomando como origen la posición inicial del camión. & (3 puntos)\\[-0.15mm]
(c) Determina la rapidez inicial $v_0$ del vehículo autónomo de apoyo y la posición $x^{*}$ del encuentro. & (12 puntos)
\end{tabularx}
\tightline

% Problema 2
\justifying
\problem{2} Un cohete de prueba se dispara verticalmente desde un pozo. El lanzamiento consta de cuatro etapas. En la primera, el cohete comienza su movimiento con una rapidez inicial de $V_{I0}=70\,\mathrm{m/s}$ a nivel del suelo. Una vez alcanzada la altura máxima en esta primera etapa, comienza la segunda etapa; el cohete enciende sus motores y acelera a $a_{II}=5\,\mathrm{m/s^2}$ hasta que alcanza una altitud de $h=1000\,\mathrm{m}$. En este punto se apagan los motores de propulsión y el cohete queda solamente bajo la acción de la gravedad en la tercera etapa. Finalmente, justo antes de volver a encontrarse con el piso, se encienden los motores para frenar uniformemente el movimiento y aterrizar a salvo, frenando de tal manera que justo al momento que hace contacto con el suelo el cohete queda en reposo. La etapa 4 inicia a una altura $H=300\,\mathrm{m}$ sobre el nivel del suelo.

\vspace{0.5mm}
\begin{tabularx}{\textwidth}{@{}>{\raggedright\arraybackslash}X>{\raggedleft\arraybackslash}p{25mm}@{}}
(a) Determina la altura máxima $y_{\max}$ que alcanza el cohete. & (10 puntos)\\[-0.15mm]
(b) Si $y_{\max}=1382{,}65\,\mathrm{m}$, determina la \underline{velocidad} con la que el cohete inicia la etapa final de aceleración, $\vect{v}_{IV}$, en su forma analítica. & (6 puntos)\\[-0.15mm]
(c) Si la rapidez con la que comienza la etapa IV es $v_{IV}=145{,}67\,\mathrm{m/s}$, determina analíticamente la aceleración ($\vect{a}_{\mathrm{final}}$) del cohete en la última etapa de aterrizaje. & (4 puntos)
\end{tabularx}
\tightline

% Problema 3
\justifying
\problem{3} Un equipo de ingenieros está delimitando el terreno para una construcción. El desplazamiento total requerido para completar el cierre perimetral no se conoce directamente y debe ser calculado a partir de las etapas previas del recorrido realizado por un dron de reconocimiento. El dron parte desde el punto de origen (estación base) y realiza los siguientes desplazamientos sucesivos. $\vect{A}$: $90\,\mathrm{m}$ en una dirección de $37^{\circ}$ al Norte del Este, $\vect{B}$: $70\,\mathrm{m}$ en una dirección de $53^{\circ}$ al Norte del Oeste y $\vect{C}$ un tercer tramo que no fue registrado por el GPS. Se sabe que, tras realizar estos tres desplazamientos, el dron se encuentra exactamente a $\vect{D}$: $40\,\mathrm{m}$ en dirección Sur respecto a la estación base.

\vspace{0.5mm}
\begin{tabularx}{\textwidth}{@{}>{\raggedright\arraybackslash}X>{\raggedleft\arraybackslash}p{25mm}@{}}
(a) Determina, de manera analítica, los desplazamientos $\vect{A}$, $\vect{B}$ y $\vect{D}$. & (6 puntos)\\[-0.15mm]
(b) Utilizando los resultados del inciso anterior, determina de manera analítica $\vect{C}$. & (8 puntos)\\[-0.15mm]
(c) Determina la magnitud y la dirección, con respecto al eje $+x$ (Este), del desplazamiento $\vect{C}$. & (6 puntos)
\end{tabularx}

\vspace{0.55mm}\hrule height 0.45pt\vspace{0.65mm}
\textbf{Ecuaciones:}
\vspace{0.15mm}
\[
\sin\theta=\frac{\mathrm{cat.\ op.}}{\mathrm{hip.}}\ ;\qquad
\cos\theta=\frac{\mathrm{cat.\ ady.}}{\mathrm{hip.}}\ ;\qquad
\tan\theta=\frac{\mathrm{cat.\ op.}}{\mathrm{cat.\ ady.}}\ ;\qquad
\vect{r}(t)=\vect{r}_0+\vect{v}_0t+\frac{1}{2}\vect{a}t^2\ ;\qquad
\vect{v}(t)=\vect{v}_0+\vect{a}t\ ;\qquad
v^2=v_0^2+2a(\Delta r)
\]
\end{document}
